\section{Conclusion}
\label{sec:conclusion}

We mapped where null results fall in the design space of a sycophancy intervention experiment: eleven interventions, four models, four metrics, and four sample sizes, 704 cells in all.
Nulls cluster heavily by evaluation metric (0.81 for a sampled answer against 0.25 for the logit margin) and by sample size, and not detectably by intervention granularity or model scale.
Most nulls are not missed behavioral effects.
Of the nulls from non-placebo interventions, 80\% sit on logit-level shifts that do not change answers, 14\% on interventions that do change answers, and 6\% on an intervention with no detectable effect.
Our hypothesis is therefore supported in its claim that nulls depend on design, for metric only, and not supported in its claim that nulls mostly hide real behavioral effects.

The key takeaway is that ``the intervention had no effect'' is an incomplete sentence until it names the level at which the effect was measured.

Several follow-ups would test how far this holds.
The most valuable is a dose-response for each intervention, to measure how many \latent units become \behav at a higher dose without breaking the model.
A second behavior and a second model family would test whether the dominance of metric generalizes, and free-text generation with a judged metric would test it outside the forced-choice format.
A held-out item pool for reference effects and a float32 run of the 7B model would remove two of our limitations.
